\section{Apparatus}\label{sec:apparatus}

This section fixes the vocabulary used in the rest of the paper.  Most
of it is imported from earlier papers of the series and is used as a
black box; we state only what the later sections need.

\subsection{Closure formation}\label{sec:apparatus:closure-formation}

Six Birds Theory treats a theory as a fixed point of an idempotent
closure operator acting on descriptions; this calculus is developed in
\FOne\ \cite{Tsiokos2026EmergenceCalculus}.  The companion papers apply
it to a mathematical problem by recording, for a \emph{substrate} (the
problem's mathematical setting), the carriers, observables, comparison
maps and audit data visible at a chosen layer.  When this typed record
is assumed to be admissibly closed under the framework's formation
discipline, we call it a \emph{formed closure}, written
\(\operatorname{FormedClosure}(S)\); the companion papers use the term
informally, and the predicate is a definition of this paper.
We call the standing assumption that the formed closures used by the
companion papers exist the \emph{SB closure assumption}.  It is a
framework assumption, not a theorem of ordinary set theory; it is not
proved here, and by itself it supplies none of the problem-specific
hypotheses of the three applications.

An \emph{accepted import} is outside mathematical content that the
closure formation allows a proof attempt to use.  It must be named and
kept apart from what the attempt produces itself.  To keep names
distinct from truth values, \Cref{sec:foreclosure} works with a set
\(\mathsf{Claim}\) of \emph{claims}, each with an interpretation
\(\llbracket c\rrbracket\) as a proposition.  The accepted imports form
a set of claims with a soundness hypothesis,
\[
  \mathcal B\subseteq\mathsf{Claim},\qquad
  c\in\mathcal B\ \Rightarrow\ \llbracket c\rrbracket .
\]
Two different claims may have equivalent interpretations, and a true
claim need not be accepted.  (A predicate on propositions would not
allow this: since equivalent propositions are equal, a sound predicate
on propositions accepts either every true proposition or none.)  The phrase
``imported as named recognition'', used for the recognition sources of
\Cref{sec:recognition}, means that a source is admitted in this way and
is not claimed to be derived.

\subsection{The Foundations III audit calculus}\label{sec:apparatus:birdint}

\FThree\ \cite{Tsiokos2026FoundationsIII} constructs a finite audited
calculus, \(\mathbf{BirdInt}^{\mathrm{aud}}_{\mathrm{fin}}\), in which
claims are recorded together with the instrument that audits them and a
verdict.  We use one of its results as motivation.  Theorem~24 of
\FThree\ (same-level self-audit failure) concerns the claim classifier
of that calculus: a well-formed complete self-soundness claim of an
instrument, judged at its own level, is not accepted; it receives
\texttt{outside\_scope}, \texttt{undefined\_circular}, or a status of
higher priority (\texttt{blocked}, \texttt{absent\_with\_record} or
\texttt{failed\_audit}).  In \Cref{sec:foreclosure} each proof attempt
\(C\) records the verdict of its same-level audit, and
\(\operatorname{SelfAuditFail}(C)\) is the statement that this recorded
verdict is \texttt{outside\_scope} or \texttt{undefined\_circular}; the
restriction to these two verdicts is an assumption of this paper, not a
consequence of Theorem~24.  The verdict
is part of the data of \(C\); it is not derived here from Theorem~24,
which explains why it is the relevant one: an attempt that can certify
its target only by auditing itself at the target's level has fallen
back on the content it was meant to supply.  The admissibility
background for these audits is that of \FTwo\
\cite{Tsiokos2026FoundationsII}.

\subsection{Adequacy residuals}\label{sec:apparatus:adequacy-residual}

The adequacy-residual companion \cite{Tsiokos2026XiCompanion} measures
how much of a target readout is left unexplained by a family of native
probes.  A finite-dimensional real or complex carrier space \(E\)
carries an \emph{audit energy} \(A\succeq0\), the cost of \(u\in E\)
being \(u^{*}Au\); the native probes form an operator \(L\) on \(E\), and
the target readout an operator \(D\), both required to vanish on
\(\ker A\) (\emph{null legality}: \(\ker A\subseteq\ker L\cap\ker D\)).  With the block currencies \(K_{FG}=FA^{\dagger}G^{*}\),
the residual is the generalized Schur complement
\[
  \Xi_{A}(D\mid L)=K_{DD}-K_{DL}K_{LL}^{\dagger}K_{LD}.
\]
The companion writes the energy as a subscript \(C\); here \(C\) is
reserved for carriers.  When the target readout is the identity and the
audit energy is the identity, \(\Xi(L):=\Xi_{I}(I\mid L)\) is the
orthogonal projection onto \(\ker L\), and
\Cref{sec:gaussian-bridge,sec:triad-bridge} use it in that form.  In
\Cref{sec:foreclosure} only one fact about the residual is used, namely
whether it vanishes.

\begin{definition}[Carrier residual]\label{def:carrier-residual}
For a carrier with audit energy \(A\), native probes \(L\) and target
readout \(D\), the \emph{carrier residual} is the adequacy residual
\(\Xi_{A}(D\mid L)\) of \cite{Tsiokos2026XiCompanion}; under null
legality it is the part of the target readout that the native probes do
not explain at the audit energy \(A\).
\end{definition}

A \emph{residual closure} is the statement \(\Xi_{A}(D\mid L)=0\).
When \Cref{sec:foreclosure} compares a residual closure with a target,
the comparison is logical: it asks whether the proposition
\(\Xi_{A}(D\mid L)=0\) is equivalent to the target.  A residual closure
does not by itself prove the target.

From the emergence companion \cite{Tsiokos2026NeedleKiller} we use the
distinction between probes that stay native to a carrier and probes that
\emph{dissolve the layer} in which the carrier operates.  A
layer-dissolving probe is not evidence that the original carrier closed
its target; its endpoint must either be recorded as outside content or
be admitted by the closure formation.

\subsection{Typed contexts}\label{sec:apparatus:typed-context}

Every claim, carrier, bridge, calibration and obligation in this paper
carries an explicit type: its source, its target, the kind of carrier
it uses, and its admissible imports.  This \emph{\TCD}, inherited from
\FTwo\ and \FThree\
\cite{Tsiokos2026FoundationsII,Tsiokos2026FoundationsIII}, forbids
replacing one of these by an informally similar one.  We write
\[
  C\leadsto X,\qquad C\models\text{\TCD}
\]
for ``the carrier \(C\) is aimed at the target \(X\)'' and ``\(C\) is
audited in a context whose source, target, carrier type and imports are
fixed''.  Two consequences are used repeatedly: a carrier aimed at one
target is not a carrier aimed at another without an explicit bridge or
re-aiming step, and a bridge from one setting to another does not by
itself give a bridge back.

\subsection{Theorems and conditional schemas}\label{sec:apparatus:opacity-guards}

The numbered statements of this paper are of two kinds.  A
\emph{theorem}, proposition or corollary is proved, from the hypotheses
it states.  A \emph{conditional schema} is a statement that is not
proved; it is displayed in its own environment, with its hypotheses, so
that later results can name it as a hypothesis.  Several predicates in
\Cref{sec:foreclosure,sec:recognition} are \emph{primitive}: they are
named but not defined in this paper, and nothing is assumed about them
beyond what a statement says explicitly.  A result about primitive
predicates establishes only the logical shape it displays.
\Cref{app:formalization} lists, for each numbered statement, what has
been formally verified.

\subsection{Glossary}\label{sec:glossary}

The following framework terms recur.  Each is used only in the sense
given here.

\begin{description}[leftmargin=1.5em,style=nextline,itemsep=2pt]
\item[Substrate.] The mathematical setting of a problem, for instance
  the zeros of the completed zeta function or the runs of machines on
  SAT instances.
\item[Formed closure; SB closure assumption.] A typed record of a
  substrate's carriers, observables, comparison maps and audit data,
  assumed admissibly closed (\Cref{sec:apparatus:closure-formation}); the
  assumption that the formed closures used by the companion papers
  exist.
\item[Carrier; residual.] A presentation through which a proof attempt
  tries to reach a target; the part of the target readout that the
  carrier's native probes leave unexplained
  (\Cref{def:carrier-residual}).
\item[Accepted import; claim.] Outside content that a closure admits; it
  is recorded as a claim, a name with an interpretation as a
  proposition (\Cref{sec:apparatus:closure-formation}).
\item[Layer-dissolving probe.] A probe whose endpoint lies outside the
  layer in which the carrier operates, so that it is not evidence that
  the carrier closed its target.
\item[Recognition source; recognition-mode landing.] A load-bearing
  structural fact that a closure names as a premise rather than
  deriving; a conditional closure built on such a premise
  (\Cref{sec:recognition}).
\item[Trace-state-only; witness-content-exposing.] The two columns of
  \Cref{def:vdiff-tso-column}: whether current observation sees only
  trace bookkeeping of a computation, or a witness datum itself.  For
  \(P\) versus \(NP\), current observation by polynomial-time machines
  is placed in the first column: it is not assumed to see a satisfying
  assignment.
\item[Calibration; functorial gate.] A typed detector of obstructions
  together with its native carrier; a named transport between
  calibrations or types whose proof is open (\Cref{sec:calibration}).
\end{description}
